\documentclass[aps,pra,twocolumn,superscriptaddress,longbibliography]{revtex4-2}

\usepackage{amsmath,amssymb,amsthm,graphicx,xcolor,bm,braket,dcolumn,booktabs}
\usepackage{hyperref}

\newtheorem{theorem}{Theorem}
\newtheorem{observation}{Observation}

% Numeric-registry citation macro (CLAUDE.md SS15): renders the literal
% only; the key is a foreign key into debug/qa/numeric_registry.py.
\newcommand{\gvq}[2]{#2}

\begin{document}

\title{The Bessel-Moment Period Algebra of the Three-Center Coulomb
Integral:\\
Modular Structure and Cosmic-Galois Placement}
\thanks{Paper 61 in the Geometric Vacuum series}

\author{J.~Loutey}
\affiliation{Independent Researcher, Kent, Washington}

\date{August 16, 2026}

\begin{abstract}
The companion Paper~59~\cite{loutey_paper59} establishes that the genuine
three-center electron-repulsion integral over $1s$ Slater orbitals, read in
momentum space under Fock's conformal projection, is a two-scale elliptic
Bessel moment: its radial content is a period of the rank-four irregular
Picard--Fuchs connection of the genus-one curve
$y^{2}=(c_{1}k^{2}+1)(c_{2}k^{2}+1)$, and its finite closed form is open.  This
paper takes that object as its starting point and develops its number-theoretic
structure.  \textbf{[MEASURED; classical]} The elliptic family is the Legendre
family over the modular curve $X(2)=\Gamma(2)\backslash\mathfrak{H}^{\ast}$,
pulled back along the exact linear map $\lambda(\tau(\rho))=1-\rho$; the
parameter integral sweeps the physical arc $\lambda<1$ of the real locus of
$X(2)$ and passes through complex-multiplication fibres---at $\tau=i$ ($\rho=\tfrac12$, inside the
physical domain) and $\tau=i\sqrt2$---where the periods are Chowla--Selberg
$\Gamma$-values, so a molecular integral carries classical $\Gamma$-value
periods across distinct fundamental discriminants.  \textbf{[OBSERVATION]} The
integrated observable lives on $X_{0}(2)$, is a $\Gamma(2)$ resurgent Lambert
series (the pure multiple modular value being its regular $D\to0$ shadow), and
carries conductor-$4$ ($\mathbb{Z}[i]$, $G=\beta(2)$) arithmetic through the
fibre theta series $K=\tfrac{\pi}{2}\theta_{3}^{2}$. That modular and
complex-multiplication structure---not the integer-relation search---places it
inside the cosmic-Galois / mixed-elliptic-motive layer the program's Tannakian
substrate reaches for;\ the guarded, decoy-calibrated searches are separately
and decisively negative for a clean low-height closed form through weight
three.
\textbf{[SYMBOLIC + MEASURED]} The connection has four master periods whose
Wronskian is the pure monomial $W_{0}D^{-2}$ and whose Lagrange bilinear
concomitant---the analogue, here for the $\Gamma(2)$ (Legendre) family, of the
Broadhurst--Roberts / Fres\'an--Sabbah--Yu quadratic period relations, whose
proved case is the Kloosterman connection---is identified in closed form as
$B=\pi\,\Omega$, with $\Omega$ the canonical integer symplectic intersection
form, so the \emph{monodromy} group lies in $\mathrm{Sp}_{4}(\mathbb{Z})$
(the differential Galois group, being a positive-dimensional algebraic
group here, lies in $\mathrm{Sp}_{4}(\mathbb{C})$).
The object is an Eisenstein / complex-multiplication Bessel-moment period,
neither a cusp-form $L$-value nor an elliptic dilogarithm; the explicit finite
reduction remains the open frontier.
\end{abstract}

\maketitle

%% ========================================================================
\section{Introduction: the object, and the tower above it}
%% ========================================================================

The genuine three-center electron-repulsion integral $(XY|XZ)$ over $1s$ Slater
orbitals is, read in momentum space under Fock's conformal projection, a
two-scale elliptic Bessel moment: a period of the rank-four irregular
Picard--Fuchs connection of the genus-one curve
$y^{2}=(c_{1}k^{2}+1)(c_{2}k^{2}+1)$, whose two scales are the
Feynman/K\"all\'en--Lehmann parameters $c_{1}=s(1-s)$ and $c_{2}=t(1-t)$.  The
construction of that object---the momentum reduction, the transcendence jump
from the genus-zero two-center class, the elliptic-dilogarithm obstruction, and
the differential-Galois irreducibility of the connection---together with its
chemistry context and full numerical validation is the companion
Paper~59~\cite{loutey_paper59}, which this paper takes as established and as its
starting point.

Here we develop the number-theoretic tower that sits above it.
Section~\ref{sec:modular} places the elliptic family as the Legendre family over
$X(2)$, exhibits its rational modulus map and its in-domain
complex-multiplication fibres, and reads the integrated observable as a
$\Gamma(2)$ resurgent Lambert series carrying conductor-$4$ arithmetic---its
cosmic-Galois placement.  Section~\ref{sec:bessel_algebra} makes that placement
quantitative through the period algebra of the master family: the Wronskian
$W_{0}D^{-2}$, the quadratic period relations, and the intersection form
$B=\pi\,\Omega$ that forces the \emph{monodromy} group into
$\mathrm{Sp}_{4}(\mathbb{Z})$, and the differential Galois group into
$\mathrm{Sp}_{4}(\mathbb{C})$.

Throughout, ``Paper~59'' denotes the companion; the equations and sections of
Paper~59 that this development rests on (the genus-one curve, the rank-four
Picard--Fuchs connection and its one-mass slice $N(D)$, the two-mass Bessel
moment, and the obstruction analysis) are pointed to in prose rather than by
cross-reference.  The tier tags (\textbf{[SYMBOLIC]}, \textbf{[MEASURED]},
\textbf{[OBSERVATION]}, \textbf{[OPEN]}) are carried over from the companion
unchanged.

\emph{Transcendence bookkeeping.}  \textbf{[OBSERVATION]} The transcendentals
below are placed against the program's taxonomy
(Papers~18,~34~\cite{loutey_paper18,loutey_paper34}), as the companion Paper~59
places its own.  Paper~18's \emph{embedding} tier carries a \emph{genus} grading, and this
object sits one rung above the \emph{two-centre} class:\ the two-center
constants are the genus-zero set $\{E_{1},\ln,\gamma_{E}\}$, while the
three-center fibre --- which this paper shares with its companion, whose own
title places it at genus one --- leaves that tower for the elliptic one.  Concretely:\ $\pi$ enters at the measure
tier of Paper~34 and appears here as the overall factor of the intersection form
$B=\pi\Omega$ and in $K=\tfrac{\pi}{2}\theta_{3}^{2}$;\ the Jacobi theta
constants $\theta_{2},\theta_{3},\theta_{4}$ and the constant
$\varpi\equiv K(\tfrac12)$ (this corpus's symbol;\ the classical
lemniscate constant is $\sqrt2\,\varpi$) are \emph{fibre} data of the Legendre family, not
projection constants;\ the $\Gamma$-values $\Gamma(1/4)^{2}$ and
$\Gamma(1/8)\Gamma(3/8)$ are Chowla--Selberg periods attached to the
complex-multiplication fibres and are therefore \emph{arithmetic} rather than
analytic in origin;\ $\zeta$ and $\beta$ appear only through the factorization
$L(\theta_{3}^{2},s)=4\zeta(s)\beta(s)$ of the $\mathbb{Z}[i]$ theta series, so
Catalan's constant $G=\beta(2)$ is conductor-$4$ \emph{family} data;\ and
$\gamma_{E}$ appears only as coordinate bookkeeping in the weight-one seed set
(\S\ref{sec:modular}), a reading this paper records as an observation
rather than a re-tiering of Paper~18's ownership of that constant.  The
load-bearing consequence is the one the searches measure:\ \emph{no} measured
conductor-$4$ content ($\beta(2)$) survives in the integrated physical value at
reachable heights, so the arithmetic lives in the family and in the $D\to0$
shadow rather than in the observable.

%% ========================================================================
\section{The modular structure and its cosmic-Galois periods}
\label{sec:modular}
%% ========================================================================

The genus-one curve of Paper~59 is not an arbitrary elliptic curve:
the family over the scale-ratio base is a \emph{modular} family, and its periods
are periods of a specific modular curve.  Forming the modular parameter from the
two periods of Paper~59, $\tau(\rho)=iK(\rho)/K(1-\rho)$, the
theta-series modular lambda satisfies \textbf{[MEASURED, precision-limited
($2.7\times10^{-51}$ at 50 digits, $1.2\times10^{-91}$ at 90); classical]}
%
\begin{equation}
  \lambda\big(\tau(\rho)\big) = 1-\rho ,
  \label{eq:lambda_rho}
\end{equation}
%
an exact, rational---indeed linear---modulus map.  The family is therefore the
Legendre family, the universal elliptic curve over the modular curve
$X(2)=\Gamma(2)\backslash\mathfrak{H}^{\ast}$, pulled back along
$\lambda=1-\rho$; it is non-isotrivial ($j(\rho)$ varies) with three cusps at
$\rho\in\{0,1,\infty\}$ matching the three cusps of $\Gamma(2)$ (the cusp
$\rho=1$ is the coincident-scale diagonal).  The periods $K(1-\rho),K(\rho)$ are
the $\Gamma(2)$ modular periods.

\emph{The modulus is a four-point cross-ratio.}  Identity~\eqref{eq:lambda_rho}
has an elementary source.  Under $u=x^{2}$ the quartic model reduces $2\!:\!1$ to
the elliptic curve $v^{2}=u(u-1)(\rho u+1-\rho)$, whose periods are
$K(\rho),K(1-\rho)$ and whose four branch points are
$\{0,1,\infty,(\rho-1)/\rho\}$.  Their six-element cross-ratio orbit is
$\{(\rho-1)/\rho,\;1/\rho,\;\rho/(\rho-1),\;\rho,\;1-\rho,\;-1/(\rho-1)\}$
\textbf{[SYMBOLIC]}, and it contains \emph{both} period arguments $\rho$ and
$1-\rho$; the Legendre value $1-\rho$ is one of its members.  The elliptic
modulus of the three-center integral is therefore the classical modulus of four
marked points on $\mathbb{P}^{1}$---the tame four-subspace ($\tilde D_{4}$)
modulus, equivalently the cross-ratio of the four regular singular points of a
rank-two Fuchsian system (the isomonodromic setting whose deformation is
Painlev\'e~VI).  One caveat fixes the reading: the four points are the branch
data of the \emph{two} spectral masses $c_{1},c_{2}$ through $\rho=c_{2}/c_{1}$
(each mass contributes a $\pm$ pair, one pair sent to $\{0,\infty\}$ and the
other to $\{1,(\rho-1)/\rho\}$ by the reduction), \emph{not} the four positions
of four scattering centers---the modulus counts the two densities, not the
number of nuclei.

Since $\rho=c_{2}/c_{1}$ ranges over all of $(0,\infty)$ as the Feynman
parameters range over $(0,1)$, the parameter integral sweeps the arc
$\lambda<1$ of the real locus of $X(2)$ --- the physical range
$\rho\in(0,\infty)$ gives $\lambda=1-\rho<1$, so the arc $\lambda>1$ is not
reached --- and passes through infinitely many of its
complex-multiplication fibres (not every one:\ those with non-real $\lambda$,
such as discriminant $-3$, are off the contour, as is $\lambda=2$;\ see
below in this section).  There the
periods are $\Gamma$-values by Chowla--Selberg \textbf{[MEASURED, $\sim\!10^{-51}$;
classical]}: at $\tau=i$ (the fibre $\rho=\tfrac12$, \emph{inside} the physical
domain), $K(\tfrac12)=\Gamma(1/4)^{2}/(4\sqrt\pi)$; at $\tau=i\sqrt2$
(discriminant $-8$), $K=(1+\sqrt2)^{1/2}\Gamma(1/8)\Gamma(3/8)/(2^{13/4}\sqrt\pi)$.
A molecular three-center integral thus has, at the fibres its own parameter
integral visits, periods that are classical $\Gamma$-values---across distinct
fundamental discriminants, so the arithmetic tracks the fibre rather than a single
coincidence.

\emph{The integrated integral is an elliptic Bessel moment.}  The full
parameter integral pulls back, via Eq.~\eqref{eq:lambda_rho}, to an iterated
integral over $X(2)$ of period-weighted forms.  The fibre, however, carries an
irregular exponential (Bessel) twist---it degenerates to the Bessel value
$K_{0}(1)$ as $\rho\to0$, which is not a modular period---so the object is
properly a \emph{Bessel moment of the Legendre family}, an irregular/exponential
period in the sense of Fres\'an--Sabbah--Yu~\cite{fresansabbahyu2023}, of which
the $\Gamma(2)$ multiple modular value in the sense of Brown~\cite{brown_mmv2014}
(an elliptic polylogarithm~\cite{brown_levin2011,beilinson_levin1994}) is the
regular $D\to0$ shadow.  A dedicated
evaluator---a fixed Gauss--Legendre tensor grid with a $\sin^{2}$ substitution on
the smooth bulk and, at the $(s,t)\to(0,0)$ corner, a Duffy angular split followed
by the radial substitution $\rho=\sigma^{2}$ ($\rho=s+t$) that renders the
$\rho^{3/2}$ corner integrand \emph{analytic} in $\sigma$ and so restores spectral
Gauss--Legendre convergence---gives the collinear value \textbf{[MEASURED,
$18$ digits confirmed]} $\gvq{t2_collinear_period}{0.3953557659017139641}\ldots$ (correct to $18$ digits;\ the
$19$th is superseded by the certified value below).  The earlier $\sim\!16$-digit
figure was not an intrinsic wall but an under-resolution of the corner and of the
one wide sub-domain flagged above (both closed here; the corner by $\rho=\sigma^{2}$,
the sub-domain by cross-validated Gauss--Legendre and $\tanh$--$\sinh$ quadrature);
the $\sim\!20$-digit ceiling was a fixed-fibre-grid artifact, not intrinsic:\ on the
decay-scaled map $k=Lu/(1-u)$, $L=1/(\sqrt{c_{1}}+\sqrt{c_{2}})$, a fixed Gauss--Legendre
fibre is \emph{spectral} to $40$-plus digits (the fibre is not the bottleneck), and a
rebuilt spectral evaluator ($\rho=\sigma^{2}$ radial + $\sin^{2}$ angular corner, kink-free
trapezoid split) independently reproduces the value to $16$ digits (the two structurally
distinct evaluators agree to $16$ digits; a naive single-axis outer-grid push produces
non-monotonic near-coincidences past that and is not trusted); the residual precision limit
is the outer-integral Gauss--Legendre rate ($\sim e^{-0.55N}$, complex-singularity-limited), so
within the $(s,t)$ frame a definitive evaluation required a different representation---supplied
below by the exact $(k,w)$ refactorization, Eq.~\eqref{eq:kw}, not by the modular route.  \textbf{[MEASURED]} A guarded
integer-relation search with a same-magnitude decoy control finds the value is
\emph{not} a low-height combination of the small disc-$4$ complex-multiplication
ring $\{1,\pi,\,K(\tfrac12)=\Gamma(\tfrac14)^{2}/4\sqrt\pi\}$ at weight $\le2$: the
relation height is far above low and is matched by the decoy, the signature of no
low-height relation---and this leg is decisive, since a weight-$\le2$ disc-$4$
basis is small enough to resolve at the reachable precision.  The wider searches,
through weight three and against the disc-$8$-inclusive ring
$\{\pi,K(\tfrac12),P_{8}\}$, are only \emph{consistent} with the negative at
present: those bases are over-determined at $\sim\!19$ digits---the spurious
low-height hits differ from one precision to the next, and are matched by the
decoy---a limitation removed below.  All of this is consistent with a genuine
(non-classical, multi-fibre) elliptic Bessel-moment period.

\textbf{[SYMBOLIC + MEASURED]} Both the precision wall and the pending
weight-three decision are closed by an exact refactorization.  Writing
$j_{0}(kb)=\int_{0}^{1}\cos(kbw)\,dw$ decouples the $b=s+t$ phase, and the
reflection symmetry $P(s,k)=P(1-s,k)$ of the collinear packet makes the inner
integral real, so the $(s,t)$ double integral collapses to the \emph{square of a
one-dimensional integral}:
\begin{equation}
T2=\frac{8}{\pi}\int_{0}^{\infty}\!\!dk\int_{0}^{1}\!\!dw\,\cos(kw)\,R(k,w)^{2},
\quad
R(k,w)=\int_{0}^{1}\!\cos\big(kw(s-\tfrac12)\big)\,P(s,k)\,ds.
\label{eq:kw}
\end{equation}
In this frame the complex off-axis $(s,t)$ singularity that produced the
$\sim\!13$--$14$-digit outer wall does not exist (the $k$-integrand is entire in
$w$;\ the $s$-integrand is analytic on $[0,1]$), no fibre is ever evaluated, and
the entire precision question is the closed-form large-$k$ tail (a Watson
expansion;\ the Gevrey structure relocates to a truncation floor $e^{-K/2}$ with
the cut $K$ freely chosen).  Six parameter-disjoint runs---two independent
parallel configurations agreeing to $1.7\times10^{-67}$, the identity itself
verified against a fully independent two-dimensional evaluation to $69$--$96$
digits---certify
\[
T2=\gvq{t2_collinear_period}{0.395355765901713964325229296804847564260563977867082108935234265469}\ldots
\]
($66$ digits;\ a \emph{decomposed} certification---identity, quadrature and tail legs
each closed independently, at higher precision than any single end-to-end pipeline;\
\texttt{tests/test\_paper59\_t2\_value.py} pins the identity, the $K^{-7}$ tail law, and the
digit-$19$ correction of the frozen anchor above, and the certified entry with its full
evidence accounting is in the reference artifact,
\texttt{benchmarks/certified\_reference/}).
At this precision the guarded, decoy-calibrated integer-relation searches become
decisive across the board:\ the corrected weight-$\le3$ ring
$\{\pi,K(\tfrac12)^{\pm1},G\}$ is \textbf{decisively negative} at height
$\le10$ ($64$ working digits, two precisions), as are disc-$4$ through weight
three and disc-$8$ through weight two at heights up to $10^{4}$--$10^{12}$, and
eight targeted Catalan probes---\emph{no measured conductor-$4$
($\beta(2)$/Catalan) content exists in the integrated observable at clean
heights}.  This sharpens the shadow reading (the finite
weight-three-requiring-$G$ element belongs to the $D\to0$ shadow, not to $T2$)
and answers the Paper~56 seam test T-2 in the negative at the clean-height
bound.  Two calibration riders:\ a decision on the dimension-twenty ring was
never available at the previously projected $\sim\!32$--$40$ digits (height
budgets scale as $10^{D/n}$ in the ring dimension $n$), and the negative
excludes \emph{clean} closed forms, not large-height ones.  The resurgent
trans-series program below is untouched---digits and closed form are separate
questions;\ the closed form remains [OPEN].

\textbf{[MEASURED]} Two refinements sharpen this.  \emph{First}, the fibre
precision wall is removable.  The edge oscillation of $j_{0}(k(s+t))$---which runs
over a $k$-range $\sim\!1/\sqrt{c_{\min}}$ as a Fock scale collapses and limited
every fixed- and decay-grid fibre quadrature to $\sim\!16$--$22$ digits---is
eliminated by an \emph{oscillation-free} evaluator:\ split the radial integral as
$\int_{0}^{K}+\int_{K}^{\infty}$ and evaluate the oscillatory tail
$\int_{K}^{\infty}\sin(kb)\,e^{-Ak}\,k^{-n}\,dk=\operatorname{Im}\!\big[z^{\,n-1}
\Gamma(1{-}n,zK)\big]$ in closed form ($z=A-ib$, $A=\sqrt{c_{1}}+\sqrt{c_{2}}$,
$b=s+t$), the tail amplitudes taken from the $1/k$-series of the
cancellation-stabilised combination $g(k)=P_{1}(k)P_{2}(k)\,e^{Ak}k^{6}$ (a single
$\Gamma(-6,zK)$ plus downward recurrence).  This reproduces the decay-grid fibre to
$40$--$45$ digits at every $(s,t)$, the corner included, and reproduces the collinear
value to its confirmed digits.  With the fibre no longer the bottleneck, the residual
precision wall of that frame was the \emph{outer} $(s,t)$-integral's corner-limited
convergence---which held the value near a $\sim\!19$-digit ceiling there until the $(k,w)$
refactorization of Eq.~\eqref{eq:kw} dissolved the question---not the fibre.  This outer wall is now characterized:\ of the
four corners only $(s,t)\to(0,0)$ is fractionally singular ($\rho^{3/2}$, absorbed by the
$\rho=\sigma^{2}$ Duffy)---the three oscillatory corners ($b=s+t\ne0$) are integer-leading (leading
order $\rho^{2}$) and carry no fractional-power subtraction \textbf{[MEASURED;
\texttt{test\_paper59\_resurgence\_corner.py}]}---so the cap is a weak \emph{complex off-axis}
singularity of $J$ (its signature a non-monotone Gauss--Legendre rate), not a missing corner map;
a third, coverage-corrected momentum-space evaluator independently reproduces the value to
$\sim\!13$--$14$ digits and then plateaus, confirming the anchor's leading digits;\ no
$(s,t)$-frame quadrature reached further, and the definitive digits were ultimately supplied by
the $(k,w)$ refactorization (Eq.~\eqref{eq:kw}), not by the modular route.  \emph{Second}, the guarded searches above used a
\emph{period-only} ring.  A length-two, weight-$\le3$ iterated integral of
weight-two Eisenstein series over $X(2)$---the shape forced by the two Feynman
integrations---generically produces also the Legendre \emph{quasiperiod}
$1/K(\tfrac12)$ (the second-kind companion, $E(\tfrac12)=\pi/4K(\tfrac12)+
K(\tfrac12)/2$) and the weight-two Eisenstein $L$-value $G=\beta(2)=L(2,\chi_{-4})$
(Catalan's constant; native to this family because the fibre period \emph{is} the
$\mathbb{Q}(i)$ theta series:\ $K(\lambda)=\tfrac{\pi}{2}\,\theta_{3}^{2}$ with
$\theta_{3}^{2}=1+\sum_{n\ge1}r_{2}(n)\,q^{n}$, $r_{2}(n)=4\sum_{d\mid n}\chi_{-4}(d)$ by
Jacobi's two-squares theorem---the weight-one, level-$4$ Eisenstein series of character
$\chi_{-4}$, with $L(\theta_{3}^{2},s)=4\,\zeta(s)\beta(s)$.  The often-quoted
$\int_{0}^{1}K(k)\,dk=2G$ holds in the \emph{modulus} convention only; in this paper's
parameter convention---which is also the physical $d\rho$ measure of the co-area
reduction below---$\int_{0}^{1}K\,dm=2$ and $\int_{0}^{1}E\,dm=\tfrac43$, both rational,
so the convention-free reason $G$ enters is the $\chi_{-4}$ theta series, not those
integrals \textbf{[MEASURED; \texttt{test\_paper59\_theta\_chi4.py}]}), neither of which
lies in $\{1,\pi,K(\tfrac12),P_{8}\}$.  The conductor-$4$ arithmetic is thus carried by
the fibre period at \emph{every} fibre---a property of the family, not of the $\tau=i$
point---and the physical Bessel twist sharpens it:\ the decay rate
$A=\sqrt{c_{1}}+\sqrt{c_{2}}=\sqrt{u}\,(1+\sqrt{\rho})$ makes
$\sqrt{\rho}=\sqrt{1-\lambda}=\theta_{4}^{2}/\theta_{3}^{2}$---a $\Gamma(4)$-level
modular function---physical, promoting the level from $\Gamma(2)$ to $\Gamma(4)$, whose
field is $\mathbb{Q}(\mu_{4})=\mathbb{Q}(i)$, with $4=|\mathrm{disc}\,\mathbb{Q}(i)|
=\mathrm{cond}\,\chi_{-4}$ \textbf{[MEASURED; \texttt{test\_paper59\_theta\_chi4.py}]}.
In the corrected ring
$\{\pi,\,K(\tfrac12),\,1/K(\tfrac12),\,G\}$ the weight-$\le2$ search is decisively
negative, and the period-only weight-$\le3$ search is negative, so the finite closed
form, should it exist, is \emph{weight-three and requires} $G$---and it is then the closed form of
the $D\to0$ \emph{shadow}, the residual question being whether the physical $D=1$ value itself
collapses onto that ring element (since answered in the negative at clean heights by the
$64$-digit decoy-calibrated search above).  This
sharpened target also recasts the earlier merely-``consistent-with'' negative as a
ring-incompleteness artifact rather than, by itself, evidence of non-classicality:\ the
non-classicality is carried separately, by the fibre's irregular twist---established
symbolically at the operator level in Paper~59
(Newton polygon, Poincar\'e rank one) and witnessed numerically below---not by the
integer-relation search.

\textbf{[OBSERVATION]} This places the three-center integral inside the
elliptic (genus-one) layer of the cosmic-Galois / mixed-elliptic-motive structure
that the Geometric Vacuum program's Tannakian substrate reaches
for~\cite{brown_cosmic2017,loutey_paper56}---the concrete elliptic period the
mixed-Tate substrate could not itself supply.  \textbf{[OPEN]} The natural next step is therefore to place it in the
Bessel-moment period algebra---to test whether it obeys a Broadhurst--Roberts
quadratic relation or reduces to a critical $L$-value of the associated modular
form~\cite{fresansabbahyu2023,zhou_quadratic2021}---of which the
explicit $\Gamma(2)$ iterated-Eisenstein
realization~\cite{broedel_eisenstein2018,manin2006} captures only the regular
$D\to0$ shadow, both of a piece with the finite closed form of
Paper~59.  We note the level is
$\Gamma(2)$---the classical \emph{universal} Legendre level, below the two-loop
sunrise's $\Gamma_{1}(6)$; landing on the Legendre family is expected for a
four-branch-point curve, so the content here is the rational modulus map and the
in-domain complex-multiplication fibres, not an exotic level.

\textbf{[OBSERVATION]} The route to that closed form is now specific.  The fibre period
is the $\Gamma(2)$ modular quantity $K(m)=\tfrac{\pi}{2}\,\theta_{3}(0,q)^{2}$ at
$\tau=iK'(m)/K(m)$ (with $\lambda(\tau)=m$, both validated to $40$ digits), so the modulus
integration pulls back to the $\tau$-plane with the weight-two Jacobian in closed form,
$d\lambda/d\tau=i\pi\,\theta_{2}^{4}\theta_{4}^{4}/\theta_{3}^{4}
=i\pi\,\lambda(1-\lambda)\theta_{3}^{4}$ (the classical modular-lambda derivative; leading
$q$-expansion $i\pi(16q-256q^{2}+2112q^{3}-\cdots)$) \textbf{[MEASURED, $\sim\!10^{-32}$
finite-difference {\rm +} symbolic; \texttt{test\_paper59\_bd\_jacobian.py}]}.  The \emph{physical}
fibre, however, carries the
$D=1$ exponential (Bessel) twist the pure period lacks---it degenerates to $K_{0}(1)$, not a
modular period, and this irregularity is measured directly at the level of the master period
itself:\ its large-$D$ asymptotic series is Gevrey-one (the optimal-truncation order grows
linearly with $D$ and the series then diverges) with Borel radius $2$, matching the dominant
singularity of the connection of Paper~59 \textbf{[MEASURED;
\texttt{test\_paper59\_resurgence\_corner.py}, with the irregular-connection witness
\texttt{test\_L4\_irregular\_at\_infinity}]}---so the integrated observable is a $\Gamma(2)$
\emph{resurgent Lambert series}
rather than a bare modular period, the pure $\Gamma(2)$ multiple modular value being its regular
$D\to0$ shadow.  This is exactly the class for which Broadhurst and Dorigoni give the
iterated-Eisenstein / Lambert-series machinery for sunrise and banana
integrals~\cite{broadhurst_dorigoni2026}.  With the Jacobian above in closed form the
length-two-to-length-one reduction is now explicit:\ a co-area change of variables (modulus
$\rho=c_{2}/c_{1}$ outer, scale $u=c_{1}$ inner) collapses the observable to a single
$X(2)$-locus integral $T2=(16/\pi)\int_{0}^{1}\Phi(\rho)\,d\rho$, folded by the exact symmetry
$\Phi(\rho)=\Phi(1/\rho)/\rho^{2}$ \textbf{[MEASURED, reduction $2.2\times10^{-9}$ vs the anchor,
symmetry $5\times10^{-29}$; \texttt{test\_paper59\_coarea\_reduction.py}]}.  The twist whose
Eichler integral this is, is the \emph{scale-integrated two-mass} fibre $\Phi(\rho)$---the
one-mass $N(D)=\int_{1}^{\infty}e^{-Dx}/\sqrt{Q}\,dx$ is only its $\rho\to0$ slice---and its cusp
sits at the diagonal $c_{1}=c_{2}$, where the curve degenerates to genus zero, so the leading
Fourier--Whittaker (log) coefficient is the genus-zero diagonal value
$A=J(\tfrac14,\tfrac14,1)/4$ (with the exact relation $b_{1}=-A$), one transcendence level below
the observable.  The cusp expansion $\Phi=\sum_{m}\!\big[b_{m}\lambda^{m}\ln\lambda
+e_{m}\lambda^{m}\big]$ is \emph{asymptotic}:\ its leading reliable terms overshoot
($68\%\!\to\!51\%$ of $T2$, non-monotone), which excludes any convergent low-order closure and
leaves the \emph{Borel--Lambert resummation} to the second cusp---the Stokes / median-summation
data one level beyond the coefficients (the BD $\sum_{n}a(n)q^{n}/(1-q^{n})$ form)---as the single
remaining specialist step.

\textbf{[SYMBOLIC + MEASURED]} The fold is itself a modular statement, and it locates the
observable's modular home one level lower than the family's.  With $c=s(1-s)$ invariant under
$s\mapsto1-s$, the full symmetry group of the Feynman square acts on the modulus
$\rho=c_{2}/c_{1}$ through exactly $\{\mathrm{id},\ \rho\mapsto1/\rho\}$;\ the physical
$s\leftrightarrow t$ density exchange is $\rho\mapsto1/\rho$, i.e.\
$\lambda\mapsto\lambda/(\lambda-1)$, the coset $T\Gamma(2)$, whose elliptic representative
$M=\bigl[\begin{smallmatrix}1&-1\\2&-1\end{smallmatrix}\bigr]$ satisfies $M^{2}=-I$ and
$M\equiv T\ (\mathrm{mod}\ 2)$, and $\langle\Gamma(2),M\rangle=\Gamma_{0}(2)$.  The folded
co-area integral $\int_{0}^{1}\Phi\,d\rho$ is therefore an integral over a fundamental domain
for the extended group:\ \emph{the observable lives on $X_{0}(2)=X(2)/\langle\iota\rangle$, not
on $X(2)$} (\texttt{tests/test\_paper59\_gamma0\_2.py}).  Two consequences, one positive and
one deflationary.  Positive:\ $\Gamma_{0}(2)$ has a unique elliptic point, of order two with
$M^{2}=-I$---disc $-4$, $j=1728$, CM by $\mathbb{Z}[i]$---so the conductor-$4$ arithmetic attaches to
the observable's modular home canonically, complementing the fibre-independent
$\theta_{3}^{2}$ leg above.  Deflationary:\ that elliptic point sits at $\lambda=2$
($\rho=-1$), the ramification point of $X(2)\to X_{0}(2)$, \emph{off} the physical contour
$\rho>0$;\ the involution's physical fixed locus is the diagonal $\rho=1$, the pure-Tate cusp;\
and $\rho\mapsto1-\rho$ is not induced by any physical symmetry, so the CM fibre
$\rho=\tfrac12$ is one member of the $s\leftrightarrow t$ orbit $\{\tfrac12,2\}$---a visited
fibre, not a pinned one.  (That a $\mathbb{Z}_{2}$ symmetry of a modular family pins a $\mathbb{Z}[i]$ point at
its elliptic fixed locus is $\mathrm{PSL}_{2}(\mathbb{Z})\cong\mathbb{Z}_{2}\ast\mathbb{Z}_{3}$ torsion---order two
forces disc $-4$, order three disc $-3$---so the arithmetic content of the involution is one
bit, not a mechanism.)

\textbf{[MEASURED]} Two literature templates fix that step precisely and, between them, locate the
one genuinely-open corner.  A single-scale reduction of exactly this kind---a Bessel moment whose
argument is a modular function, equated to a weight-two modular form and thence to an Eichler
integral of a higher-weight form---is Zhou's~\cite{zhou_wick2018} (e.g.\ $\int_{0}^{\infty}\!
J_{0}(\alpha(z)t)\,I_{0}(t)K_{0}(t)^{3}\,t\,dt=\tfrac{\pi^{2}}{16}Z_{6,3}(z)$ at level
$\Gamma_{0}(6)$, and the weight-four Eichler representation of $\int I_{0}K_{0}^{4}$), whose special
values close to complex-multiplication $\Gamma$-values or critical $L$-values.  The resummation of
the resulting cusp series across the two Fricke-dual cusps is Broadhurst--Dorigoni's, both in the
divisor-sum form they give for sunrise/banana integrals~\cite{broadhurst_dorigoni2026} and---closer
to our setup---in the character-twisted precursor~\cite{broadhurst_dorigoni_characters2025}, whose
transseries relates the value at one cusp to the same Lambert family at the Fricke-dual argument
$\tau\to-1/(r_{1}r_{2}\tau)$, median-resummed with algebraic (character) Stokes data; the general
engine---a Borel tower whose Stokes constants are $L$-function coefficients, resummed to a quantum
modular form---is the modular-resurgence paradigm of Fantini--Rella~\cite{fantini_rella2024} and the
Eichler-integral median resummation of McSpirit--Rolen~\cite{mcspirit_rolen2025}.  Our object sits
one axis beyond each:\ the fibre of Paper~59 is a \emph{two-scale} ($a_{1}\neq a_{2}$)
Bessel moment where Zhou's are single-scale, and its cusp coefficients are Bessel--Fourier--Whittaker
coefficients rather than the Dirichlet divisor sums those resummations twist by.  The distinction is
concrete already at the leading coefficient:\ the diagonal value
$A=J(\tfrac14,\tfrac14,1)/4=0.0790540321168\ldots$ is itself a single-scale Bessel moment---by
the $K_{0}$ identity of Paper~59 and its $t$-derivative $\int_{0}^{\infty}\!\cos(kb)\,e^{-t\sqrt{ck^{2}+1}}\,dk=
2t\,K_{1}(\sqrt{t^{2}+4b^{2}})/\sqrt{t^{2}+4b^{2}}$, $A$ reduces (exactly) to a
$\sigma$-integral of $K_{1}$ at argument $\sqrt{t^{2}+4\sigma^{2}}$---and a guarded, decoy-controlled
integer-relation search finds \emph{no} low-height closure for it in the weight-one two-center ring
$\{\pi,e^{-a},K_{0},K_{1},E_{1}\}$ at the natural arguments $a\in\{2,2\sqrt2\}$
\textbf{[MEASURED; \texttt{test\_routeC\_momentum.py::test\_paper59\_diagonal\_A}]}.  So even the leading cusp datum is an
irreducible Bessel moment, not an elementary or divisor-sum quantity---consistent with a resurgent
Lambert series whose individual coefficients need not be elementary for the resummed value to be a
clean $\Gamma(2)$ period.  \textbf{[OPEN]} The exact remaining object is therefore the weight-two
Eisenstein Eichler integral over $X(2)$ against this two-scale Bessel twist---the two-scale
generalization of the single-scale Bessel-moment theory of
Fres\'an--Sabbah--Yu~\cite{fresansabbahyu2023}, itself open---covered by none of these templates
individually but pinned by all of them jointly.

\textbf{[MEASURED]} The second-cusp resurgence itself is now characterized in closed form.  The
scale-integrated twist's one-mass expansion ($c_{2}\!\to\!0$, one scale collapsing) is Gevrey $(2n)!$-type, and
its dominant Borel singularity---the Stokes-data \emph{location}---is
$z^{*}=-(\sqrt{c_{1}}\mp ib)^{2}=(b^{2}-c_{1})\pm 2i\sqrt{c_{1}}\,b$, with $|z^{*}|=c_{1}+b^{2}$ and
trans-series action $\sqrt{z^{*}}=b+i\sqrt{c_{1}}$ (the source separation $b$ and the scale
$\sqrt{c_{1}}$ fused into one complex action).  It is a $(z^{*}-w)^{3/2}$ branch
(exponent $p=-\tfrac52$, measured cleanly) with amplitude in \emph{closed form},
$2|A|=(c_{1}+b^{2})^{3/2}/(4\sqrt{\pi}\,\sqrt{c_{1}}\,b)$ (algebraic $\times\,1/\sqrt\pi$, the
$1/\sqrt\pi$ from the $(1+\varepsilon)^{-5/2}$ boundary branch)---derived from the moment tail rate
$\sqrt{c_{1}}-ib$ and validated at three parameter points (fit-to-closed-form ratio
$0.988$--$0.989$, a uniform finite-order deficit converging to $1$).  This is the two-scale analog of the algebraic-Stokes closure of $N(D)$ in
Paper~59, and the median-Borel resummation across the (complex, off-contour)
singularity pair is unambiguous.  \textbf{[OPEN]} The residual step to the finite value is the
integration of this algebraic-Stokes trans-series over the Feynman family (the inner scale) and the
modulus---the concrete form the two-scale Broadhurst--Dorigoni resummation takes here.  Brute
high-precision quadrature is not a route \emph{in the $(s,t)$ or co-area frames}
(complex-singularity- and fibre-cost-limited there);\ the $(k,w)$ refactorization
Eq.~\eqref{eq:kw} removes exactly that obstruction and supplies the \emph{digits}, while the
finite \emph{closed form} remains this resummation question.

\textbf{[MEASURED]} \emph{The Stokes location is a complexified light cone.} In position
space each packet is analytic in $b^{2}$ with branch points at the endpoint of its
source-time superposition, where the Euclidean interval closes, $b^{2}+p_{i}^{2}=0$; the
$j_{0}$ smearing moves the singularity to the endpoint pinch
$\beta^{*}=\pm|\bm W|\pm i p_{1}$, so the Borel location above is the squared complexified
interval, $z^{*}=(\beta^{*})^{2}$, $|z^{*}|=p_{1}^{2}+|\bm W|^{2}$, of which
$-(\sqrt{c_{1}}\mp ib)^{2}$ is the $D=1$ specialisation $p_{1}=\sqrt{c_{1}}$; the trans-series
action $a_{2}\sqrt{z^{*}}$ is the heavy mass times that interval. Moment growth of the
position-space fibre confirms modulus and phase at eight $(c,|\bm W|,D)$ points, including
$D\neq1$, where the interval form and the $D$-independent form differ by $35\%$
\textbf{[MEASURED, $0.6$--$5\%$ from twelve coefficients]}. Two consequences. The median
resummation is unambiguous because a Euclidean interval never vanishes at real separation:
the absence of a Stokes ambiguity is Euclidean positivity. And the $(s,t)\to(0,0)$ corner is
the coincidence (ultraviolet) point, where the interval closes and the Borel singularities
collide at the origin---which is why that corner, and only that corner, obstructs the family
integration. There the masses diverge at fixed ratio, so the modulus survives and the Duffy
angle \emph{is} the modulus; the leading behaviour is the ultraviolet power count
$J=\sigma^{3}A(a)+O(\sigma^{5})$, $\sigma^{3}=m_{1}^{-2}m_{2}^{-2}m$, so the $\rho^{3/2}$
non-analyticity is an odd mass dimension and carries no $\pi$. At the three corners with
$|\bm W|\neq0$ the oscillation survives, a Dirichlet limit supplies one further power and an
explicit $\pi$, and the integer-leading behaviour is now closed:
$J=\sigma^{4}\,\frac{\pi}{2}G(1)^{2}\alpha(1-\alpha)/|\bm W|+O(\sigma^{5})$ with $G(1)=7/e$
\textbf{[MEASURED, $10^{-4}$ with $O(\sigma^{2})$ approach, four corner/angle combinations;\
\texttt{test\_paper59\_euclidean\_dictionary.py}]}.

\textbf{[OBSERVATION]} Read against the two-layer taxonomy, the object separates on two
kinematic axes: the mass split turns on the elliptic (genus-one) level and the Euclidean-time
displacement turns on the irregular level---the curve carries no $D$, and the large-$D$ series
is Gevrey-one at every $\rho$, including the genus-zero slice $\rho=0$ where the fibre is
$K_{0}(D)$---while the modulus then also grades the irregular data through the Borel radius
$\min(2,\,a_{2}/a_{1})$, the smallest difference of the four exponential rates. Every $\pi$
in the object comes from the momentum-space measure or from integration over the momentum
half-line; the non-compact Euclidean-time direction contributes no $\pi$ but the exponential
and its Stokes data, consistent with the compactification reading of
Paper~35~\cite{loutey_paper35}.  (In Euclidean signature the labels are structural, not
metric:\ $D$ is the shift each propagator carries, $b$ the shared coordinate that is
integrated out;\ the temporal reading is offered as a reading.)

\textbf{[OBSERVATION]} The algebraic-Stokes closures of this paper sit inside a pattern that now
spans the corpus's two-electron integral classes.  Four objects have had their resurgent data
computed:\ the Paper-18 exchange-class seed $e^{a}E_{1}(a)$ (Borel transform $1/(1+\zeta)$, one
simple pole, Stokes constant $2\pi i\cdot1$);\ the Paper-58 hybrid two-center class
$\{E_{1},\ln\}$ (all Borel singularities at $\pm a_{d}$, the far-center exponent;\ Stokes
constants $2\pi i\times\mathbb{Q}(\sqrt d\,)$, extracted blind from large-order coefficients to
$35$--$81$ digits and matched by exact integer relations;\ and the sector constants cancel
pairwise, so the physical integral is \emph{cut-free} in $R$---the median resummation \emph{is}
the closed form, \texttt{tests/test\_paper59\_resurgent\_skeleton.py});\ the one-mass fibre
$N(D)$ of Paper~59 (algebraic over $\mathbb{Q}(\rho)$);\ and the second-cusp
triple above (algebraic $\times\,1/\sqrt\pi$).  In all four the connection data is algebraic
over the parameter field up to a $\pi$-power that is itself fixed by the Borel singularity type
(simple pole $\to2\pi i$;\ square-root branch $\to\pi^{0}$;\ three-halves branch
$\to1/\sqrt\pi$)---nothing in the Stokes data is free.  What varies is the \emph{boundary
period}:\ absent for the seed;\ $\ln\Lambda$ with $\Lambda\in\mathbb{Q}$ the multiplicative
cross-ratio of the four Borel actions---i.e.\ \emph{forced by the skeleton}---for the hybrid
class;\ the elliptic $K(1-\rho)$, \emph{not} determined by the Stokes data, for $N(D)$;\ and the
open $\Gamma(2)$ object for $T2$.  The forced$\to$free transition of the boundary period thus
sits between weight one and genus one---the same seam at which the closed-form program of this
paper opened.  The named next test---the exchange class, whose $R$-dependent logarithm makes
its Borel singularity logarithmic rather than polar---has since been run and the pattern stands
at \textbf{five for five}:\ the exchange kernel collapses to a single-sector reduced normal
form whose Borel transform is the closed-form, charge-neutral sum
$\mathfrak{B}(\xi)=\tfrac{1}{2ab}\sum_{j}q_{j}(\xi-\xi_{j})\ln(\xi-\xi_{j})$ over the four
rational positions, with \emph{integer} Stokes charges $q_{j}\in\{-1,+1,+1,-1\}$ (no radical
needed, unlike the hybrid's $\sqrt d$);\ the boundary log is forced even more tightly than the
hybrid's cross-ratio---its argument is $\kappa=\prod_{j}\lambda_{j}^{q_{j}}=2ab/(a{+}b)$, the
exponents \emph{being} the Stokes charges;\ and Euler's $\gamma$ never exists in the Borel
plane at all:\ the $R$-dependent log moves one Borel singularity onto the sector's own origin
(pole $\to u\ln u$ branch), and $\gamma$ is exactly the bookkeeping cost of writing that origin
singularity in the $R$ variable---which finally tags the $\gamma$ of the weight-one seed set as
coordinate bookkeeping, not a projection transcendental.  The object is Borel-summable with no
lateral ambiguity yet genuinely multivalued in complex $R$ (the price of the $R$-dependent log,
costing no new transcendental), where the hybrid class was single-valued
(\texttt{tests/test\_paper59\_resurgent\_skeleton.py}).  The $\pi$-power law refines
accordingly:\ an \emph{integral} local Borel exponent (pole \emph{or} log branch) fixes $2\pi
i$;\ a half-integral one fixes $\pi^{0}$ or $1/\sqrt\pi$.  (Read against the two-layer taxonomy
this is a candidate refinement---a resurgent-skeleton grading \emph{within} the transcendental
layer---recorded as an observation, now on five objects;\ the residual open leg is the
heteronuclear infinite $\tau$-sum.)

%% ========================================================================
\section{The Bessel-moment period algebra: master family and quadratic
relations}
\label{sec:bessel_algebra}
%% ========================================================================

Section~\ref{sec:modular} places the integrated observable in the Bessel-moment
class; we now make that placement quantitative by exhibiting the period algebra
it inhabits.  The rank-four irregular connection of Paper~59 has four master
periods, the Laplace transforms
%
\begin{equation}
s_{c}(D)=\int_{\gamma_{c}}\frac{e^{-Dx}\,dx}{\sqrt{Q(x)}},
\qquad Q(x)=(x^{2}-1)(\rho x^{2}+1-\rho),
\end{equation}
%
over the four independent cycles $\gamma_{c}$ joining the branch points
$\{\pm1,\pm i\sqrt{(1-\rho)/\rho}\}$: the physical moment $s_{K}$ (the cut
$[1,\infty)$, which is $N(D)$ itself, $\sim e^{-D}$), its companion $s_{I}$
($[-1,1]$, $\sim e^{+D}$), the imaginary-axis period $s_{J}$
($\sim\cos\omega D$), and the sector $s_{Y}$ carrying the $D\ln D$
non-analyticity.  \textbf{[MEASURED]} The three regular masters satisfy
the Picard--Fuchs equation of Paper~59 to $\sim\!10^{-15}$ across $\rho$ and $D$ (each
annihilated: $s_{K}\!\sim\!10^{-18}$, $s_{I}\!\sim\!10^{-17}$,
$s_{J}\!\sim\!10^{-15}$).

\emph{The Wronskian.}  \textbf{[SYMBOLIC]} The operator of Paper~59 is formally
self-adjoint, so its differential Galois group lies
in $\mathrm{Sp}_{4}$, and Abel's identity---with the sub-leading coefficient
ratio $p_{3}/p_{4}=2/D$---fixes the Wronskian of the master matrix exactly as a
pure monomial,
%
\begin{equation}
W(D)=W_{0}\,D^{-2}:
\end{equation}
%
the $D$-dependence of the determinant is the pure monomial $D^{-2}$, the
elliptic analogue of the classical $W[K_{0},I_{0}](D)\propto 1/D$ and the
two-scale counterpart of the Wro\'nskian factorizations behind the
Broadhurst--Mellit determinant formulae~\cite{zhou_wronskian2018}.
%
\emph{Correction, 2026-09-07}
\texttt{[retracted 2026-09-07: p61-w0-transcendence-cancels]}.\ Earlier versions
added that ``the transcendence of the individual masters cancels in their
determinant''. That is \emph{false}, and the value below is what establishes
it:\ Abel's identity fixes only the
$D$-dependence, and the constant $W_{0}$ is \emph{never computed} in this paper
or in its backing test, which normalises it to $1$ by construction. It is now
computed, from the same branch-point data that fixes $B$ below. With
$Q(x)=(x^{2}-1)(\rho x^{2}+1-\rho)$ and branch points $\{\pm1,\pm i\omega\}$,
$\omega^{2}=(1-\rho)/\rho$:\ $\sum_{c}x_{c}=0$ cancels the exponentials;\
$\prod_{c}D^{-1/2}=D^{-2}$ reproduces Abel's exponent by an independent route;\
$\prod_{c}Q'(x_{c})=-16\omega^{2}$, so $\prod_{c}A_{c}=\pi^{2}/(4i\omega)$;\ and
the Vandermonde $\prod_{a<b}(x_{b}-x_{a})=4i\omega/\rho^{2}$ cancels the
$\omega$, leaving
%
\begin{equation}
W_{0}\;=\;\frac{\pi^{2}}{\rho^{2}}
\qquad \textbf{[SYMBOLIC]},
\label{eq:W0}
\end{equation}
%
up to the thimble ordering/orientation sign this paper already riders for $B$
(the sign is the branch of $\sqrt{1-\rho}/\sqrt{\rho-1}=-i$ on $0<\rho<1$).
So $\pi^{2}$ does \emph{not} cancel --- it is squared --- and $W_{0}$ is
$\rho$-dependent, which is what makes the withdrawn sentence false rather than
merely unsupported. The withdrawal is a net gain:\ an exact value replaced a
false cancellation.

\emph{Quadratic period relations.}  \textbf{[MEASURED, 25 digits]}
Self-adjointness makes the Lagrange bilinear concomitant $B[s_{a},s_{b}]$
constant in $D$; evaluated on the masters it is moreover independent of the
modulus $\rho$ and lands on rational multiples of $\pi$,
%
\begin{equation}
B[s_{K},s_{I}]=-\pi,\qquad B[s_{K},s_{J}]=0,\qquad B[s_{I},s_{J}]=2\pi,
\end{equation}
%
verified to $25$ digits at $D=1$ across the tested $\rho$, and $D$-independent
by the self-adjoint Lagrange identity.
These are the Broadhurst--Roberts / Fresán--Sabbah--Yu
quadratic relations between Bessel
moments~\cite{fresansabbahyu2023,zhou_quadratic2021}, here for the
$\Gamma(2)$ (Legendre) family---the elliptic lift of $W[K_{0},I_{0}]=1/D$.
Structurally the concomitant is the de~Rham--Betti intersection pairing of the
thimble cycles times the branch-point half-residue $\pi$ (the square-root
$(-1)$-monodromy of $1/\sqrt{Q}$ at each simple branch point, the elliptic lift of
the $1/\pi$ normalizing $I_{0}$ and of $W[K_{0},I_{0}]=1/D$), so $B=\pi\,J$ with $J$
an integer intersection form.

\textbf{[SYMBOLIC]} \emph{The intersection form, identified.}  We close the
identification of $J$ from the branch-point asymptotics---no numerics enter
in this leg.  Near a
simple branch point $Q\sim Q'(x_{c})(x-x_{c})$, so each master has the leading Laplace
behaviour $s_{c}(D)\sim e^{-x_{c}D}\sqrt{\pi/(Q'(x_{c})\,D)}$ (the half-order integral
$\int_{0}e^{-Dt}t^{-1/2}\,dt=\Gamma(\tfrac12)/\sqrt D$); since the concomitant is exactly
constant in $D$, its value equals its $D\to\infty$ limit, which these leading terms
capture (the subleading $O(1/D)$ corrections vanish in the limit).  In the
Lefschetz-thimble basis $\{\gamma_{+1},\gamma_{-1},\gamma_{+i\omega},\gamma_{-i\omega}\}$
---one thimble per branch point, the basis in which the monodromy $M_{0}$ of
Paper~59 is the exact integer matrix---two facts follow.
\emph{Block-diagonality:} $s_{a}s_{b}\sim e^{-(x_{a}+x_{b})D}$, so a nonzero
\emph{constant} forces $x_{a}+x_{b}=0$; the branch points pair as $\{+1,-1\}$ (the real
$K_{0}/I_{0}$ sector) and $\{+i\omega,-i\omega\}$ (the imaginary $J_{0}/Y_{0}$ sector),
and every cross-sector pairing ($x_{a}+x_{b}\neq0$) vanishes identically.
\emph{Plane-equality and the unit $\pi$:} on a within-sector pair the leading concomitant
is $B[s_{x},s_{-x}]=2x\,(-2\rho x^{2}+2\rho-1)\,A_{x}A_{-x}$ with
$A_{x}A_{-x}=\pi/\sqrt{Q'(x)Q'(-x)}$, and the exact identity
$Q'(x)Q'(-x)=-4x^{2}(2\rho x^{2}-2\rho+1)^{2}$ cancels the algebraic factor completely,
leaving
%
\begin{equation}
B[s_{x},s_{-x}]^{2}=-\pi^{2}\qquad\text{independently of }x\text{ and }\rho,
\label{eq:planepi}
\end{equation}
%
so \emph{both} symplectic planes carry the same pairing, of magnitude exactly $\pi$ (the
$\pi=\Gamma(\tfrac12)^{2}$ of the two branch-point half-integrals).  Hence
%
\begin{equation}
J=\Omega:=\begin{pmatrix}0&1\\-1&0\end{pmatrix}\oplus
\begin{pmatrix}0&1\\-1&0\end{pmatrix},
\qquad B=\pi\,\Omega ,
\end{equation}
%
the canonical block symplectic form; in the period-cut basis this is the real
$-\pi,0,2\pi$ measured above (a numerical thimble normalization instead reads
$B=\pi(\rho i)\Omega$, the same $\Omega$).  This $\Omega$
is \emph{forced}, not fitted: the antisymmetric forms preserved by the monodromy,
$M_{0}^{\mathsf T}J M_{0}=J$, constitute a four-parameter family, and imposing the
concomitant's vanishing cross-sector pairings selects \emph{exactly}
$\mathbb{Z}\,\Omega$.  It is unimodular ($\det\Omega=1$, hence nondegenerate), and
$M_{0}^{\mathsf T}\Omega\,M_{0}=\Omega$ holds identically over $\mathbb{Z}$, so the
\emph{monodromy} group of the Picard--Fuchs equation of Paper~59 lies in the
\emph{integral} symplectic group
$\mathrm{Sp}(\Omega,\mathbb{Z})=\mathrm{Sp}_{4}(\mathbb{Z})$---the concrete,
integral sharpening of the self-adjointness that placed the Galois group in
$\mathrm{Sp}_{4}(\mathbb{C})$.
%
\emph{Correction, 2026-09-07.}\ Earlier versions of this paper said the
\emph{differential Galois group} lies in $\mathrm{Sp}_{4}(\mathbb{Z})$. That is
false:\ $\mathrm{Sp}_{4}(\mathbb{Z})$ is discrete, so a Zariski-closed subgroup
of $\mathrm{GL}_{4}(\mathbb{C})$ contained in it would be finite, forcing every
solution to be algebraic---contradicting the irregularity of the connection and
the exponential torus $(\mathbb{C}^{*})^{2}$ that Paper~59 establishes inside
the local Galois group. The integrality is a statement about the Betti lattice,
i.e.\ the monodromy representation, which is what $M_{0}$ is;\ Paper~59 states
the Galois containment correctly as $\mathrm{Sp}_{4}(\mathbb{C})$. Nothing else
changes:\ $B=\pi\,\Omega$, the forcing of $\Omega$, and the integral structure
all stand.
\emph{Tier scope.}\ The asymptotic identification above is numerics-free, but
every clause invoking $M_{0}$---the forcing argument and the $\mathrm{Sp}_{4}(\mathbb{Z})$
containment---inherits that matrix's \textbf{[SYMBOLIC + MEASURED]} status:\ $M_{0}$
is an exact integer certificate rather than a floating-point fit, but it was
obtained by computation, not derived in closed form.  The
three physical period cuts $\{[1,\infty),[-1,1],i[-\omega,\omega]\}$ realize only a
rank-two sub-block of $\pi\Omega$ (the measured $-\pi,0,2\pi$); it is the fourth,
$Y_{0}$-sector master---the one carrying the $D\ln D$---that completes the pairing to
the full nondegenerate rank-four form.  The symbolic identification of the
intersection form is thus complete: $B=\pi\,\Omega$, the Broadhurst--Roberts /
Fresán--Sabbah--Yu quadratic relations realized for the $\Gamma(2)$ family, backed by
\texttt{tests/test\_routeC\_momentum.py} (\texttt{test\_intersection\_form\_pi\_and\_planes\_are\_symbolic}
for Eq.~\eqref{eq:planepi} and the block structure,
\texttt{test\_intersection\_form\_is\_forced\_canonical\_symplectic} for the
$M_{0}$-forcing, \texttt{test\_intersection\_form\_period\_cut\_values} for the
$-\pi,0,2\pi$ values, and
\texttt{test\_paper59\_bessel\_moment\_algebra.py::test\_paper59\_period\_pairing\_minus\_pi\_0\_2pi}
witnessing them to the quoted $25$ digits).

\emph{Not a cusp-form $L$-value.}  \textbf{[OBSERVATION]} Because
$M_{\bullet}(\Gamma(2))=\mathbb{C}[\theta_{2}^{4},\theta_{4}^{4}]$ is free on two
weight-two generators, $\dim S_{k}(\Gamma(2))$ vanishes for $k=2,4$ and first
reaches $1$ at $k=6$ (the cusp form $\theta_{2}^{4}\theta_{3}^{4}\theta_{4}^{4}$).
The observable's transcendental weight is at most three (length $\le2$ over
$X(2)$, from the two Feynman integrations), whereas a weight-six cusp form's
critical $L$-values sit at motivic weight five; and the period pairing just
exhibited is Eisenstein-flavoured (pure $\pi$-rationals), not cuspidal.  The
integrated three-center integral is therefore \emph{not} a critical $L$-value of
a $\Gamma(2)$ cusp form: it is an Eisenstein / complex-multiplication period,
built from $\pi$ and the CM-fibre $\Gamma$-values of Sec.~\ref{sec:modular}
($K(\tfrac12)=\Gamma(\tfrac14)^{2}/4\sqrt\pi$, the discriminant-$8$ period, and
their kin).  \textbf{[MEASURED]} A guarded integer-relation
search (with a same-magnitude decoy control) confirms the value is not a
low-height combination of the small weight-$\le2$ disc-$4$ Eisenstein/CM ring,
nor of the small weight-$\le1$ disc-$8$ ring $\{1,\pi,\varpi,P_{8}\}$ ($P_{8}$
the discriminant-$8$ period)---both \emph{decisive} at the conservatively
cross-validated $\sim\!16$ digits: two structurally independent evaluators (the
spectral corner of Sec.~\ref{sec:modular} and a decay-scaled-fibre quadrature)
agree to $16$ digits, the corner evaluator extending single-method to
$\sim\!19$; on the decisive rings the relation height is far above low and tracks
the decoy at every tested precision.  The wider searches---through weight three
in disc-$4$, and the disc-$8$ ring at weight two---remained only \emph{consistent}
with the negative: those dimension-ten bases are over-determined at this
precision (target and decoy both turn up small but precision-\emph{unstable}
relations, changing from one precision to the next).  The decision has since
been made at $64$ digits via the $(k,w)$ refactorization
(Eq.~\eqref{eq:kw}):\ \emph{decisively negative} through weight three,
including the Catalan-completed corrected ring, at height $\le10$---so the
finite closed form, if it exists at all, either has height $>10$ in that ring
or lives outside it.

This narrows the open problem sharply: the finite closed form, should it exist,
is a reduction of the observable to $\pi$, the CM-fibre $\Gamma$-values, their
Legendre \emph{quasiperiods}, and the weight-two Eisenstein $L$-value
$G=\beta(2)$ (Sec.~\ref{sec:modular}), within the Bessel-moment period
algebra---a two-scale, shifted member of the Broadhurst--BBBG
class~\cite{bbbg2008,broadhurst2008} evaluated over $X(2)$---and is neither an
elliptic dilogarithm nor a cusp-form $L$-value.  \textbf{[OBSERVATION]} Because the corrected-ring searches are now decisively negative
through weight three at clean heights ($64$ digits, height $\le10$), any such finite
combination either belongs to the $D\to0$ \emph{shadow}---where the shadow reading of Sec.~\ref{sec:modular} places the
weight-three $G$-requiring element---or sits at height $>10$ in that ring;\ the physical
value itself carries no measured conductor-$4$ content at reachable heights.

\begin{thebibliography}{99}

\bibitem{bbbg2008}
D.~H.~Bailey, J.~M.~Borwein, D.~Broadhurst, and M.~L.~Glasser, ``Elliptic
integral evaluations of Bessel moments and applications,''
\textit{J.\ Phys.\ A}\ \textbf{41}, 205203 (2008); arXiv:0801.0891.

\bibitem{broadhurst2008}
D.~Broadhurst, ``Elliptic integral evaluation of a Bessel moment by contour
integration of a lattice Green function,'' arXiv:0801.4813 (2008).

\bibitem{broadhurst_dorigoni2026}
D.~Broadhurst and D.~Dorigoni, ``Resurgent Lambert series from Feynman and
beyond,'' PoS \textbf{LL2026}, 020 (2026); arXiv:2607.14020.

\bibitem{broadhurst_dorigoni_characters2025}
D.~Broadhurst and D.~Dorigoni, ``Resurgent Lambert series with characters,''
arXiv:2507.21352 (2025).

\bibitem{fantini_rella2024}
V.~Fantini and C.~Rella, ``Modular resurgent structures,'' arXiv:2404.11550
(2024).

\bibitem{mcspirit_rolen2025}
E.~McSpirit and L.~Rolen, ``Quantum modular forms and resurgence,''
arXiv:2505.00799 (2025).

% NOTE (2026-09-07): no longer cited. It was in the quadratic-relations
% bundle, which was wrong -- this paper is L-series of Kloosterman moments
% and determinants of Feynman integrals, not quadratic relations. Kept as
% background rather than deleted silently; PI to keep or cut.
\bibitem{broadhurst2016}
D.~Broadhurst, ``Feynman integrals, $L$-series and Kloosterman moments,''
\textit{Commun.\ Number Theory Phys.}\ \textbf{10}, 527--569 (2016);
arXiv:1604.03057.

\bibitem{fresansabbahyu2023}
J.~Fres\'an, C.~Sabbah, and J.-D.~Yu, ``Quadratic relations between Bessel
moments,'' \textit{Algebra Number Theory}\ \textbf{17}, 541--602 (2023);
arXiv:2006.02702.

\bibitem{brown_mmv2014}
F.~Brown, ``Multiple modular values and the relative completion of the
fundamental group of $M_{1,1}$,'' arXiv:1407.5167 (2014).

\bibitem{brown_cosmic2017}
F.~Brown, ``Feynman amplitudes, coaction principle, and cosmic Galois group,''
\textit{Commun.\ Number Theory Phys.}\ \textbf{11}, 453--556 (2017);
arXiv:1512.06409.

\bibitem{brown_levin2011}
F.~Brown and A.~Levin, ``Multiple elliptic polylogarithms,'' arXiv:1110.6917
(2011).

\bibitem{beilinson_levin1994}
A.~Beilinson and A.~Levin, ``The elliptic polylogarithm,'' in \textit{Motives}
(Seattle, 1991), Proc.\ Symp.\ Pure Math.\ \textbf{55}, Part 2 (American
Mathematical Society, Providence, 1994), pp.\ 123--190.

\bibitem{manin2006}
Yu.~I.~Manin, ``Iterated integrals of modular forms and noncommutative modular
symbols,'' in \textit{Algebraic Geometry and Number Theory}, Progr.\ Math.\
\textbf{253} (Birkh\"auser, Boston, 2006), pp.\ 565--597; arXiv:math/0502576.

\bibitem{broedel_eisenstein2018}
J.~Broedel, C.~Duhr, F.~Dulat, B.~Penante, and L.~Tancredi, ``Elliptic symbol
calculus:\ from elliptic polylogarithms to iterated integrals of Eisenstein
series,'' \textit{J.\ High Energy Phys.}\ \textbf{08}, 014 (2018); arXiv:1803.10256.

\bibitem{zhou_wick2018}
Y.~Zhou, ``Wick rotations, Eichler integrals, and multiloop Feynman diagrams,''
\textit{Commun.\ Number Theory Phys.}\ \textbf{12}, 127--192 (2018);
arXiv:1706.08308.

\bibitem{zhou_quadratic2021}
Y.~Zhou, ``Wro\'nskian algebra and Broadhurst--Roberts quadratic
relations,'' \textit{Commun.\ Number Theory Phys.}\ \textbf{15}(4),
651--741 (2021); arXiv:2012.03523. \emph{The quadratic relations proper;
distinct from the Broadhurst--Mellit \emph{determinant} formulae of
\cite{zhou_wronskian2018}, with which they were conflated until
2026-09-07.}

\bibitem{zhou_wronskian2018}
Y.~Zhou, ``Wro\'nskian factorizations and Broadhurst--Mellit determinant
formulae,'' \textit{Commun.\ Number Theory Phys.}\ \textbf{12}, 355--407 (2018);
arXiv:1711.01829.

\bibitem{loutey_paper18}
J.~Loutey, ``Spectral-Geometric Exchange Constants: Projecting Discrete
Spectra onto Continuous Manifolds,'' Paper~18 in this series.

\bibitem{loutey_paper34}
J.~Loutey, ``GeoVac as a Two-Layer Framework: Projection Taxonomy and the
Empirical Matches Catalog,'' Paper~34 in this series.

\bibitem{loutey_paper35}
J.~Loutey, ``Time as Projection: The Klein-Gordon Spectrum on $S^3\times\mathbb{R}$
and the Algebraic\,/\,Observation Split,'' Paper~35 in this series.

\bibitem{loutey_paper56}
J.~Loutey, ``The Tannakian substrate of \textsc{GeoVac},''
Paper~56 in this series.

\bibitem{loutey_paper59}
J.~Loutey, ``The Three-Center Electron-Repulsion Integral Is an Elliptic
Bessel Moment: Genus One at the Third Center in Momentum-Space Slater
Theory,'' Paper~59 in this series.

\end{thebibliography}

\end{document}
